\chapter*{IMO301 ve PDR209 İzleme Rotaları}
\markboth{IMO301 ve PDR209 İzleme Rotaları}{IMO301 ve PDR209 İzleme Rotaları}
\addcontentsline{toc}{chapter}{IMO301 ve PDR209 İzleme Rotaları}

Birleşik kitap ortak bir bölüm sırası kullanır; ancak iki dersin vurgu ve
çalışma derinliği farklıdır. Aşağıdaki tablo, bölümlerin izlencelerdeki
yerini gösterir. ``Temel'' ilgili ders için haftalık ana okumayı, ``destek''
ise seçilmiş alt başlık veya uygulamaları belirtir.

\small
\begin{longtable}{@{}p{0.07\textwidth}p{0.49\textwidth}p{0.16\textwidth}p{0.16\textwidth}@{}}
\toprule
\textbf{Bölüm} & \textbf{Konu} & \textbf{IMO301} & \textbf{PDR209} \\
\midrule
\endfirsthead
\toprule
\textbf{Bölüm} & \textbf{Konu} & \textbf{IMO301} & \textbf{PDR209} \\
\midrule
\endhead
1 & İstatistiksel düşünme, veri ve araştırma süreci & Temel & Temel \\
2 & Grafiksel ve sayısal betimleme & Temel & Temel \\
3 & Eksik veri ve veri kalitesi & Temel & Temel \\
4 & Örnekleme yöntemleri ve yanlılık & Temel & Temel \\
5 & Örnekleme dağılımları & Temel & Temel \\
6 & Normal dağılım, standart puanlar ve MLT & Temel & Temel \\
7 & Nokta tahmini & Temel & Temel \\
8 & Güven aralıkları & Temel & Temel \\
9 & Bootstrap ve rastgeleleştirme & Destek & Destek \\
10 & Hipotez testleri, hata, güç ve etki & Temel & Temel \\
11 & Tek, bağımsız ve eşleştirilmiş $t$ testleri & Temel & Temel \\
12 & ANOVA ve grup karşılaştırmaları & Destek & Temel \\
13 & Kategorik veri ve ki-kare & Temel & Destek \\
14 & Korelasyon & Temel & Temel \\
15 & Basit doğrusal regresyon & Temel & Temel \\
16 & Parametrik olmayan yöntemler & Destek & Temel \\
17 & Ölçek puanlarının güvenirliği & Destek & Temel \\
18 & Çoklu doğrusal regresyon & İleri okuma & Destek \\
19 & Bütünleştirici veri analizi laboratuvarı & Uygulama & Uygulama \\
20 & Genel değerlendirme ve yöntem seçimi & Tekrar & Tekrar \\
\bottomrule
\end{longtable}
\normalsize

IMO301 rotasında örnekleme dağılımları, Merkezi Limit Teoremi, standart hata
ve çıkarım mantığına; PDR209 rotasında etki büyüklüğü, ANOVA, varsayım
kontrolü, parametrik olmayan alternatifler ve bilimsel raporlamaya daha çok
zaman ayrılması önerilir. Bölümlerdeki PDR vakaları farklı alanlardan gelen
okuyucular için bağlamsal uygulama; matematiksel derinleşmeler ise yöntemin
neden çalıştığını görmek için ortak kaynak niteliğindedir.